\documentclass[11pt,a4paper]{article}
\usepackage[T1]{fontenc}
\usepackage[utf8]{inputenc}
\usepackage[main=italian,provide=*]{babel}
\usepackage{amsmath,amssymb}
\usepackage{mathtools}
\usepackage{xcolor}
\usepackage[a4paper,margin=2.7cm]{geometry}
\usepackage{caption}
\usepackage{booktabs}
\usepackage{enumitem}
\usepackage{placeins}
\usepackage[colorlinks=true,linkcolor=black,citecolor=black,urlcolor=blue!55!black]{hyperref}
\newcommand{\ket}[1]{\lvert #1 \rangle}
\newcommand{\bra}[1]{\langle #1 \rvert}

\title{Guida d'uso --- pacchetto riproducibile\\Parte 2 (rumore), trimero ad anello}
\author{Tesi triennale, Samuele --- Relatore: Prof.\ Alessandro Chiesa}
\date{}

\begin{document}
\sloppy
\maketitle

\tableofcontents

\section{Struttura del pacchetto}

\begin{verbatim}
codice/            moduli .py (libreria, mai eseguiti direttamente)
dati/              output .npz degli script di generazione dati
figure/            figure finali, PDF+PNG
genera_figure/     uno script per figura
01_..., 10_...     script di dati/validazione, numerati nell'ordine d'uso
esegui_tutto.py    orchestratore: dati -> validazioni -> figure
\end{verbatim}

Comando singolo per rigenerare tutto da zero:
\begin{verbatim}
python3 esegui_tutto.py
\end{verbatim}
Richiede: \texttt{qiskit}, \texttt{qiskit-aer}, \texttt{numpy}, \texttt{scipy}, \texttt{matplotlib}.

\section{Dipendenza incrociata con Parte 1, dichiarata}

Cinque moduli in \texttt{codice/} sono copie bit-per-bit di file del
Project che appartengono concettualmente a \textbf{Parte 1} (sistema
chiuso), non a Parte 2: \texttt{vqe\_w2q6\_\allowbreak trimero\_anello.py},
\texttt{trimer\_ring\_exact.py}, \texttt{trotter\_\allowbreak trimero\_anello.py},
\texttt{circuito\_\allowbreak correlazioni\_\allowbreak trimero\_anello.py},
\texttt{noise\_model\_dimero.py} (quest'ultimo dal dimero, non
dall'anello). Riusati qui invece di duplicarne la logica --- vedi i
docstring dei moduli di Parte 2 che li importano per il dettaglio di
cosa viene preso da ciascuno.

\section{Le due onde di dipendenza reale}

\textbf{Onda 1} (nessun ingresso): solo \texttt{01}. \textbf{Onda 2}
(dipendono tutti da \texttt{01}, indipendenti fra loro):
\texttt{02}--\texttt{10}. \textbf{Onda 3} (dipendono dai \texttt{.npz}
prodotti nell'Onda 2): gli script in \texttt{genera\_figure/}.

\section{Nota: evoluzione del metodo di lettura (readout)}
\label{sec:nota-readout}

Gli script \texttt{01}--\texttt{07} (e i moduli che importano) applicano
l'errore di lettura \textbf{analiticamente}:
$\langle Z\rangle_\text{readout}=(1-2p)\langle Z\rangle_\text{ideale}$ ---
corretto, ma mai verificato contro una misura vera in questo pacchetto.
Gli script \texttt{08}--\texttt{10} (Sez.~\ref{sec:script-estensioni})
usano invece un \texttt{ReadoutError} vero agganciato al
\texttt{NoiseModel}, con misura a shot finiti campionata da
\texttt{AerSimulator} --- il metodo poi effettivamente citato nella
documentazione dei risultati finale. Le due strade sono state verificate
coincidere entro l'errore statistico atteso (script \texttt{08},
controllo 1) prima di scartare la formula analitica come riferimento
primario.

\section{Script 01--07}

\subsection{01\_genera\_parametri\_vqe.py}
\textbf{Ingresso}: nessuno (calcola tutto da zero: Hamiltoniana +
ottimizzazione multistart). \textbf{Uscita}:
\texttt{dati/w2q6\_params\_optimal.npz} (campi: \texttt{params},
\texttt{E\_vqe}, \texttt{E\_exact}, \texttt{fidelity}). \textbf{Chi lo
consuma}: tutti gli script \texttt{02}--\texttt{07}.
\begin{verbatim}
python3 01_genera_parametri_vqe.py
\end{verbatim}

\subsection{02\_valida\_modello\_rumore.py}
\textbf{Ingresso}: \texttt{dati/w2q6\_params\_optimal.npz}.
\textbf{Uscita}: nessun file, solo resoconto a schermo (tre controlli
sul modello di rumore, Stadio 1). \textbf{Chi lo consuma}: nessuno.
\begin{verbatim}
python3 02_valida_modello_rumore.py
\end{verbatim}

\subsection{03\_valida\_vqe\_dm\_rumoroso.py}
\textbf{Ingresso}: \texttt{dati/w2q6\_params\_optimal.npz}.
\textbf{Uscita}: \texttt{dati/scan\_eps2q\_\allowbreak vqe\_dm\_trimero\_anello.npz}
(Stadio 2). \textbf{Chi lo consuma}:
\texttt{genera\_figure/\allowbreak genera\_scan\_eps2q\_\allowbreak trimero\_anello.py}.
\begin{verbatim}
python3 03_valida_vqe_dm_rumoroso.py
\end{verbatim}

\subsection{04\_valida\_trotter\_rumoroso.py}
\textbf{Ingresso}: \texttt{dati/w2q6\_params\_optimal.npz}.
\textbf{Uscita}:
\texttt{dati/scan\_N\_trotter\_\allowbreak rumoroso\_trimero\_anello.npz} (Stadio 3,
entrambi gli scenari A/B, griglia $N$ fine --- vedi nota sotto).
\textbf{Chi lo consuma}:
\texttt{genera\_figure/\allowbreak genera\_trotter\_\allowbreak rumoroso\_trimero\_anello.py}.
\begin{verbatim}
python3 04_valida_trotter_rumoroso.py
\end{verbatim}

\subsection{05\_valida\_correlatori\_rumorosi.py}
\textbf{Ingresso}: \texttt{dati/w2q6\_params\_optimal.npz}.
\textbf{Uscita}: \texttt{dati/scan\_N\_\allowbreak correlatore\_trimero\_anello.npz}
(Stadio 4). \textbf{Chi lo consuma}:
\texttt{genera\_figure/\allowbreak genera\_correlatore\_\allowbreak rumoroso\_trimero\_anello.py}.
\begin{verbatim}
python3 05_valida_correlatori_rumorosi.py
\end{verbatim}

\subsection{06\_genera\_scan\_parametri.py}
\textbf{Ingresso}: \texttt{dati/w2q6\_params\_optimal.npz}.
\textbf{Uscita}:
\texttt{dati/scan\_parametri\_\allowbreak rumore\_trimero\_anello.npz} (Stadio 5,
entrambi gli scenari). \textbf{Chi lo consuma}:
\texttt{genera\_figure/\allowbreak genera\_Nstar\_vs\_eps\_\allowbreak trimero\_anello.py}.
\begin{verbatim}
python3 06_genera_scan_parametri.py
\end{verbatim}

\subsection{07\_valida\_scan\_parametri.py}
\textbf{Ingresso}: \texttt{dati/w2q6\_params\_optimal.npz}.
\textbf{Uscita}: nessun file, solo resoconto a schermo (regressione,
monotonia, invarianza da readout). \textbf{Chi lo consuma}: nessuno.
\begin{verbatim}
python3 07_valida_scan_parametri.py
\end{verbatim}

\section{Script 08--10 (estensioni)}
\label{sec:script-estensioni}

Coprono le due estensioni facoltative (VQE noise-aware, readout
asimmetrico) e il passaggio al readout a shot finiti --- assenti dalla
prima versione di questo pacchetto, aggiunti per allinearlo alla
documentazione dei risultati finale. Vedi Sez.~\ref{sec:nota-readout}
per il perché del cambio di metodo.

\subsection{08\_valida\_correlatori\_shots.py}
\textbf{Ingresso}: \texttt{dati/w2q6\_params\_optimal.npz}.
\textbf{Uscita}: \texttt{dati/scan\_N\_\allowbreak correlatore\_shots\_trimero\_anello.npz}
(Stadio 4bis). \textbf{Chi lo consuma}:
\texttt{genera\_figure/\allowbreak genera\_correlatore\_\allowbreak shots\_trimero\_anello.py}.
Verifica: coerenza con la formula analitica (Stadio 4) entro l'errore
statistico; $N^*=3$; convenzione dell'ancilla (qubit 3, primo carattere).
\begin{verbatim}
python3 08_valida_correlatori_shots.py
\end{verbatim}

\subsection{09\_valida\_vqe\_noise\_aware.py}
\textbf{Ingresso}: \texttt{dati/w2q6\_params\_optimal.npz}.
\textbf{Uscita}: \texttt{dati/vqe\_noise\_aware\_\allowbreak validato.npz} (Stadio 2bis).
\textbf{Chi lo consuma}:
\texttt{genera\_figure/\allowbreak genera\_vqe\_\allowbreak noise\_aware\_trimero\_anello.py}.
Verifica: l'artefatto del transpilatore (conteggio CX deve restare 6 a
struttura fissa); riproducibilità dei valori di riferimento
($\Delta E\approx-4.16\times10^{-3}$, $\Delta F\approx-1.58\times10^{-2}$);
il segno atteso ($E_\text{noise-aware}\leq E_\text{riuso}$).
\begin{verbatim}
python3 09_valida_vqe_noise_aware.py
\end{verbatim}

\subsection{10\_valida\_readout\_asimmetrico.py}
\textbf{Ingresso}: \texttt{dati/w2q6\_params\_optimal.npz}.
\textbf{Uscita}: \texttt{dati/readout\_\allowbreak asimmetrico\_validato.npz} (Stadio 6bis).
\textbf{Chi lo consuma}:
\texttt{genera\_figure/\allowbreak genera\_readout\_\allowbreak asimmetrico\_trimero\_anello.py}.
Verifica: il caso simmetrico si riduce al risultato dello Stadio 4bis;
$N^*=3$ invariante su un rapporto $p_{10}/p_{01}$ da $1\times$ a $50\times$.
\begin{verbatim}
python3 10_valida_readout_asimmetrico.py
\end{verbatim}

\section{Script genera\_figure/}

\subsection{genera\_depolarizzazione\_bloch.py}
\textbf{Ingresso}: nessuno (formula analitica $1-\lambda$).
\textbf{Uscita}: \texttt{figure/fig\_\allowbreak depolarizzazione\_bloch.pdf/.png}.
\begin{verbatim}
python3 genera_figure/genera_depolarizzazione_bloch.py
\end{verbatim}

\subsection{genera\_scan\_eps2q\_trimero\_anello.py}
\textbf{Ingresso}: \texttt{dati/scan\_eps2q\_\allowbreak vqe\_dm\_trimero\_anello.npz}.
\textbf{Uscita}: \texttt{figure/fig\_scan\_eps2q\_\allowbreak trimero\_anello.pdf/.png}.
\begin{verbatim}
python3 genera_figure/genera_scan_eps2q_trimero_anello.py
\end{verbatim}

\subsection{genera\_trotter\_rumoroso\_trimero\_anello.py}
\textbf{Ingresso}:
\texttt{dati/scan\_N\_trotter\_\allowbreak rumoroso\_trimero\_anello.npz}.
\textbf{Uscita}:
\texttt{figure/fig\_trotter\_\allowbreak rumoroso\_trimero\_anello.pdf/.png}.
\begin{verbatim}
python3 genera_figure/genera_trotter_rumoroso_trimero_anello.py
\end{verbatim}

\subsection{genera\_correlatore\_rumoroso\_trimero\_anello.py}
\textbf{Ingresso}: \texttt{dati/scan\_N\_\allowbreak correlatore\_trimero\_anello.npz}.
\textbf{Uscita}:
\texttt{figure/fig\_correlatore\_\allowbreak rumoroso\_trimero\_anello.pdf/.png}.
\begin{verbatim}
python3 genera_figure/genera_correlatore_rumoroso_trimero_anello.py
\end{verbatim}

\subsection{genera\_Nstar\_vs\_eps\_trimero\_anello.py}
\textbf{Ingresso}: \texttt{dati/scan\_parametri\_\allowbreak rumore\_trimero\_anello.npz}.
\textbf{Uscita}: \texttt{figure/fig\_Nstar\_vs\_eps\_\allowbreak trimero\_anello.pdf/.png}.
\begin{verbatim}
python3 genera_figure/genera_Nstar_vs_eps_trimero_anello.py
\end{verbatim}

\subsection{genera\_correlatore\_shots\_trimero\_anello.py}
\textbf{Ingresso}: \texttt{dati/scan\_N\_\allowbreak correlatore\_shots\_trimero\_anello.npz}.
\textbf{Uscita}: \texttt{figure/fig\_correlatore\_\allowbreak shots\_trimero\_anello.pdf/.png}.
\begin{verbatim}
python3 genera_figure/genera_correlatore_shots_trimero_anello.py
\end{verbatim}

\subsection{genera\_vqe\_noise\_aware\_trimero\_anello.py}
\textbf{Ingresso}: \texttt{dati/vqe\_noise\_aware\_\allowbreak validato.npz}.
\textbf{Uscita}: \texttt{figure/fig\_vqe\_\allowbreak noise\_aware\_trimero\_anello.pdf/.png}.
\begin{verbatim}
python3 genera_figure/genera_vqe_noise_aware_trimero_anello.py
\end{verbatim}

\subsection{genera\_readout\_asimmetrico\_trimero\_anello.py}
\textbf{Ingresso}: \texttt{dati/readout\_\allowbreak asimmetrico\_validato.npz}.
\textbf{Uscita}: \texttt{figure/fig\_readout\_\allowbreak asimmetrico\_trimero\_anello.pdf/.png}.
\begin{verbatim}
python3 genera_figure/genera_readout_asimmetrico_trimero_anello.py
\end{verbatim}

\section{Nota: correzione della griglia \texorpdfstring{$N$}{N} (Stadio 3)}

Una prima versione dello script 04 usava una griglia $N$ troppo grezza
($\{1,2,5,10,15,\ldots\}$), che saltava $N=3$ e $N=4$ e riportava
$N^*_A=2$ invece del vero $3$. Corretto con una griglia continua
$N=1,\ldots,20$ (poi diradata) --- lo script in questo pacchetto usa già
la versione corretta. Vedi
\texttt{risultati\_scan\_parametri\_rumore\_trimero\_anello.tex},
Sez.\ ``Correzione'', per i dettagli.

\section{Sequenza completa (copia-incolla)}

\begin{verbatim}
python3 01_genera_parametri_vqe.py
python3 02_valida_modello_rumore.py
python3 03_valida_vqe_dm_rumoroso.py
python3 04_valida_trotter_rumoroso.py
python3 05_valida_correlatori_rumorosi.py
python3 06_genera_scan_parametri.py
python3 07_valida_scan_parametri.py
python3 08_valida_correlatori_shots.py
python3 09_valida_vqe_noise_aware.py
python3 10_valida_readout_asimmetrico.py
python3 genera_figure/genera_depolarizzazione_bloch.py
python3 genera_figure/genera_scan_eps2q_trimero_anello.py
python3 genera_figure/genera_trotter_rumoroso_trimero_anello.py
python3 genera_figure/genera_correlatore_rumoroso_trimero_anello.py
python3 genera_figure/genera_Nstar_vs_eps_trimero_anello.py
python3 genera_figure/genera_correlatore_shots_trimero_anello.py
python3 genera_figure/genera_vqe_noise_aware_trimero_anello.py
python3 genera_figure/genera_readout_asimmetrico_trimero_anello.py
\end{verbatim}
Equivalente, in un solo comando: \texttt{python3 esegui\_tutto.py}.

\section{Verifica end-to-end effettuata}

\texttt{dati/} e \texttt{figure/} svuotate, \texttt{esegui\_tutto.py}
rieseguito da zero: \texttt{exit code 0}, nessun \texttt{FALLITO} n\'e
\texttt{Traceback} nel log. Numeri rigenerati confrontati con quelli nei
documenti di risultati: $E_\text{VQE}=-5.534370017393404$,
$N^*_A=3$ ($F=0.700335$), $N^*_B=9$ ($F=0.381896$), $N^*$
correlatore (formula analitica)$=3$ ($|C|=0.505$), $N^*$ correlatore
(shot finiti)$=3$ ($|C|=0.502$, scarto dalla formula analitica coerente
con l'errore statistico atteso), VQE noise-aware
($\Delta E=-4.16\times10^{-3}$, $\Delta F=-1.58\times10^{-2}$), $N^*$
invariante da $1\times$ a $50\times$ nel readout asimmetrico --- tutti
coincidenti.

\end{document}
